% --- WAT language definition and style ---
\lstdefinelanguage{WAT}{
  morekeywords={
    module,type,func,param,result,export,import,local,global,data,
    i32,i64,f32,f64,v128,const,set,get,load,store,unreachable,
    block,loop,if,else,br,br_if,call,offset,align
  },
  sensitive=true,
  alsoletter={.,_},
  comment=[l]{;;}
}

\lstdefinestyle{watpaper}{
  language=WAT,
  basicstyle=\ttfamily\footnotesize,
  keywordstyle=\color{blue!80!black}\bfseries,
  commentstyle=\itshape\color{green!50!black},
  stringstyle=\color{red!70!black},
  numbers=left,
  numberstyle=\tiny\color{black},
  stepnumber=1,
  numbersep=2pt,             % reduce space between numbers and code
  frame=lines,               % only top and bottom borders
  rulecolor=\color{black},
  framesep=6pt,
  xleftmargin=0pt,           % remove extra indentation
  framexleftmargin=0pt,      % align code with numbers
  showstringspaces=false,
  keepspaces=true,
  columns=fullflexible,
  breaklines=true,
  breakatwhitespace=false
}

% --- Example WAT code listing ---
\begin{lstlisting}[style=watpaper,caption={Example (Wamr) (Crash Analysis Pattern \Fix{cite})},label={example:ex7}]
(module
  (memory 65536 65536)
    ...
    i16x8.extmul_high_i8x16_u
    i8x16.extract_lane_s 1
    v128.load16_splat offset=30 align=1
    v128.store offset=691 align=1
  )
)
\end{lstlisting}

    